\chapter{API reference}
\label{ch:api}

The reference follows the groups of \code{kernel/rtos.h}. Every entry gives the
prototype, the behaviour, the parameters, the calling context, and which contracts are
\emph{checked} (always) or \emph{asserted} (\code{RTOS\_CFG\_ASSERT}). Unless an entry
says otherwise, a \code{timeout} is in ticks: \code{RTOS\_NO\_WAIT} polls,
\code{RTOS\_FOREVER} waits for good, any other value is at most
\code{RTOS\_MAX\_TIMEOUT} (checked). Functions that may wake a higher task run it before
they return; their \code{FromIsr} variants defer the switch to the end of the interrupt.
The generated reference (\code{make api} in nornDocs) has the same text per function
with the remaining details.

\clearpage
\section{Constants and types}

\begin{table}[htbp]
\centering\small
\begin{tabularx}{\linewidth}{@{}>{\raggedright\arraybackslash}p{0.38\linewidth}X@{}}
\toprule
Name & Meaning \\
\midrule
\code{RTOS\_VERSION}           & version string of the kernel \\
\code{RTOS\_IDLE\_TASK}        & 0, the idle task \\
\code{RTOS\_NUM\_TASKS}        & \code{RTOS\_CFG\_NUM\_TASKS + 1} \\
\code{RTOS\_NO\_TASK}          & 0xFF, ``no task'' \\
\code{RTOS\_STACK\_MIN / DEFAULT / LARGE} & the port's stack sizes (table~\ref{tab:stacksizes}) \\
\code{RTOS\_FOREVER}, \code{RTOS\_NO\_WAIT} & timeouts: never, poll \\
\code{RTOS\_MAX\_TIMEOUT}      & the longest finite timeout, half the tick range \\
\code{RTOS\_MS(ms)}            & milliseconds to ticks, at compile time \\
\code{RTOS\_WAIT\_ANY / ALL}   & options of \code{rtos\_waitEvents}, \code{rtos\_evgroupWait} \\
\code{RTOS\_WAIT\_CLEAR}       & \code{rtos\_evgroupWait}: take the satisfying bits out of the group \\
\code{RTOS\_NOTIFY\_SET / OR / INC} & operations of \code{rtos\_notify} \\
\code{RTOS\_SEM\_INIT(n)}, \code{RTOS\_SEM\_INIT\_BINARY(n)}, \code{RTOS\_SEM\_INIT\_MAX(n, max)} & static semaphore initialisers \\
\code{RTOS\_MUTEX\_INIT(ceiling)} & static mutex initialiser \\
\code{RTOS\_QUEUE\_INIT(storage, itemSize, capacity)} & static queue initialiser \\
\code{RTOS\_MBOX\_INIT(storage, itemSize)} & static mailbox initialiser \\
\code{RTOS\_POOL\_INIT(storage, blockSize, numBlocks)} & static pool initialiser; \code{rtos\_poolInit} is still required \\
\code{RTOS\_EVGROUP\_INIT(bits)} & static event group initialiser \\
\code{RTOS\_RUNTIME\_HZ}, \code{RTOS\_RUNTIME\_TO\_MS(t)} & run time counter frequency and conversion \\
\code{rtos\_tick\_t}, \code{rtos\_evt\_t}, \code{rtos\_taskmask\_t}, \code{rtos\_runtime\_t} & the tick, event flag, task mask and run time types; widths from the configuration and the port \\
\code{rtos\_taskFunc\_t}, \code{rtos\_timerFunc\_t} & \code{void f(void *arg)} \\
\code{rtos\_sem\_t}, \code{rtos\_mutex\_t}, \code{rtos\_queue\_t}, \code{rtos\_mbox\_t}, \code{rtos\_pool\_t}, \code{rtos\_timer\_t}, \code{rtos\_evgroup\_t} & the kernel objects; treat as opaque except for the statistics fields named in chapter~\ref{ch:robust} \\
\code{rtos\_taskInfo\_t}, \code{rtos\_panicInfo\_t}, \code{rtos\_heapStats\_t} & diagnostic records \\
\code{RTOS\_TASK\_*}, \code{RTOS\_WAIT\_*}, \code{RTOS\_RESET\_*}, \code{RTOS\_TRACE\_*} & task states, wait kinds, reset causes, trace events \\
\bottomrule
\end{tabularx}
\caption{Constants and types.}
\end{table}

% -----------------------------------------------------------------------------------------
\section{Kernel life cycle}

\apihead{rtos\_init}
\begin{lstlisting}
#define rtos_init()   rtos_initChecked(RTOS_CONFIG_SIGNATURE)
\end{lstlisting}
Resets all kernel data: every task slot, the ready mask, the wake-up list and the tick
count, and takes the panic record of the previous boot out of the memory that survives a
reset. The macro passes the configuration signature; a kernel compiled with another
configuration stops at once with \code{RTOS\_PANIC\_CONFIG}.
\apipre{Called from \code{main()}, before any other kernel function, once.}

\apihead{rtos\_createTask}
\begin{lstlisting}
void rtos_createTask(uint8_t prio, rtos_taskFunc_t taskFunc, void *arg,
                     uint8_t *stackArea, size_t stackSize);
\end{lstlisting}
Creates the task at priority \code{prio} on the caller's stack. The task is ready at once
and runs at the earliest at \code{rtos\_start}. The stack start is aligned up and the
size rounded down to the port's alignment (1 byte on the AVR, 16 on the ESP32), which
may reduce the usable size; with \code{RTOS\_CFG\_STACK\_CHECK} the area is filled with
the pattern.
\begin{apiparams}
\item[prio] task priority and id, 1 to \code{RTOS\_CFG\_NUM\_TASKS}; the slot must be free.
\item[taskFunc] the task's function; should never return (section~\ref{sec:lifecycle}).
\item[arg] passed to the function unchanged.
\item[stackArea, stackSize] caller allocated memory, valid for as long as the system
  runs; at least \code{RTOS\_STACK\_MIN} after alignment.
\end{apiparams}
\apipre{After \code{rtos\_init}, before \code{rtos\_start}. Creating a task while the
scheduler runs is not supported yet.}
\apichecked{\code{prio} out of range, a NULL function or stack, a slot in use, a stack
below the minimum.}

\apihead{rtos\_start}
\begin{lstlisting}
void rtos_start(void);
\end{lstlisting}
Creates the timer task if configured, starts the tick, switches to the highest ready
task and never returns. The calling context becomes the idle task.

% -----------------------------------------------------------------------------------------
\section{Time and delays}

\apihead{rtos\_getTick}
\begin{lstlisting}
rtos_tick_t rtos_getTick(void);
\end{lstlisting}
\apireturns{the ticks since \code{rtos\_start}. The counter wraps (16 bits: 65.5 s at
1 kHz; 32 bits: 49.7 days); compare as differences.}
\apicontext{any.}

\apihead{rtos\_delay}
\begin{lstlisting}
void rtos_delay(rtos_tick_t ticks);
\end{lstlisting}
Blocks the calling task for at least \code{ticks} ticks; the real delay is between
\code{ticks} and \code{ticks + 1} periods. A delay of 0 waits for the next tick. Any
value is accepted: a delay above \code{RTOS\_MAX\_TIMEOUT} is done as a sequence of
waits, \code{RTOS\_FOREVER} blocks for good.
\apicontext{task; not the idle task (checked).}

\apihead{rtos\_waitPeriodic}
\begin{lstlisting}
bool rtos_waitPeriodic(rtos_tick_t period);
\end{lstlisting}
Blocks until the task's next period boundary; each boundary is the previous one plus
\code{period}, so the task runs drift free. If the boundary has passed, the call returns
at once, counts an overrun and computes the next boundary from now. The first boundary
is the task's start plus \code{period}. The period may change from call to call.
\begin{apiparams}
\item[period] 1 to \code{RTOS\_MAX\_TIMEOUT}.
\end{apiparams}
\apireturns{\code{true} on an overrun, \code{false} after a normal wait.}
\apicontext{task; not the idle task.}
\apichecked{a period of 0 or above the limit.}

% -----------------------------------------------------------------------------------------
\section{Task event flags}

\apihead{rtos\_waitEvents}
\begin{lstlisting}
rtos_evt_t rtos_waitEvents(rtos_evt_t mask, uint8_t opt, rtos_tick_t timeout);
\end{lstlisting}
Waits until events of \code{mask} are pending for the calling task. Events posted
earlier stay pending until a wait consumes them. The events that satisfy the wait are
cleared and returned; pending events outside the mask are kept.
\begin{apiparams}
\item[mask] the events to wait for; not 0 (asserted).
\item[opt] \code{RTOS\_WAIT\_ANY} or \code{RTOS\_WAIT\_ALL}.
\end{apiparams}
\apireturns{the consumed events (with \code{RTOS\_WAIT\_ALL}: the mask); 0 on timeout.}
\apicontext{task; the idle task may only poll.}

\apihead{rtos\_postEvents, rtos\_postEventsFromIsr}
\begin{lstlisting}
void rtos_postEvents(uint8_t taskId, rtos_evt_t events);
void rtos_postEventsFromIsr(uint8_t taskId, rtos_evt_t events);
\end{lstlisting}
ORs the events into the task's pending events; a task whose wait is thereby satisfied
becomes ready.
\apichecked{\code{taskId} 1 to \code{RTOS\_CFG\_NUM\_TASKS}.}

\apihead{rtos\_broadcastEvents}
\begin{lstlisting}
void rtos_broadcastEvents(rtos_evt_t events);
\end{lstlisting}
Posts the events to every created task (not the idle task).
\apicontext{task.}

% -----------------------------------------------------------------------------------------
\section{Counting semaphores}

\apihead{rtos\_semInit, rtos\_semInitMax}
\begin{lstlisting}
void rtos_semInit(rtos_sem_t *pSem, uint8_t initialCount);
void rtos_semInitMax(rtos_sem_t *pSem, uint8_t initialCount, uint8_t maxCount);
\end{lstlisting}
The run time alternatives to the static initialisers. Without a limit the count is
bounded by its type, 255; with \code{maxCount} a give at the limit is dropped and counted
as lost (1 makes a binary semaphore).
\apipre{no task waits on the semaphore.}
\apiasserted{\code{initialCount} at most \code{maxCount}, \code{maxCount} at least 1.}

\apihead{rtos\_semTake}
\begin{lstlisting}
bool rtos_semTake(rtos_sem_t *pSem, rtos_tick_t timeout);
\end{lstlisting}
Takes one token, waiting up to \code{timeout} for it.
\apireturns{\code{true} if a token was taken, \code{false} on timeout.}
\apicontext{task; the idle task may only poll.}

\apihead{rtos\_semGive, rtos\_semGiveFromIsr}
\begin{lstlisting}
void rtos_semGive(rtos_sem_t *pSem);
void rtos_semGiveFromIsr(rtos_sem_t *pSem);
\end{lstlisting}
Gives one token. If tasks wait, the highest gets it directly and the count does not
change; otherwise the count goes up or, at the limit, the give is lost.

% -----------------------------------------------------------------------------------------
\section{Ceiling mutexes}

\apihead{rtos\_mutexInit}
\begin{lstlisting}
void rtos_mutexInit(rtos_mutex_t *pMutex, uint8_t ceilingPrio);
\end{lstlisting}
The run time alternative to \code{RTOS\_MUTEX\_INIT}. \code{ceilingPrio} is the highest
priority of any task that will ever lock the mutex.
\apipre{the mutex is not locked.}

\apihead{rtos\_mutexLock}
\begin{lstlisting}
void rtos_mutexLock(rtos_mutex_t *pMutex);
\end{lstlisting}
Raises the caller's effective priority to the ceiling and takes the mutex. With a correct
ceiling the call never blocks. If the mutex is held anyway (a ceiling too low, or an
owner that blocked while holding), the caller waits without a timeout and the unlock
hands the mutex over to the highest waiter.
\apicontext{task; not the idle task (asserted).}
\apiasserted{the caller does not already hold the mutex (not recursive); the caller's
priority is at most the ceiling.}
\apinote{Panics}{a wait that would close a cycle of mutex waits is
\code{RTOS\_PANIC\_DEADLOCK} (the caller, the owner).}

\apihead{rtos\_mutexUnlock}
\begin{lstlisting}
void rtos_mutexUnlock(rtos_mutex_t *pMutex);
\end{lstlisting}
Drops the caller's effective priority back to what it was before the lock. A waiting
task becomes the owner at once, at the ceiling. Either change may let another task run
before the call returns.
\apicontext{task.}
\apiasserted{the caller owns the mutex; mutexes are unlocked in the reverse order of
locking.}

% -----------------------------------------------------------------------------------------
\section{Message queues}

\apihead{rtos\_queueInit}
\begin{lstlisting}
void rtos_queueInit(rtos_queue_t *pQ, uint8_t *storage, uint8_t itemSize, uint8_t capacity);
\end{lstlisting}
The run time alternative to \code{RTOS\_QUEUE\_INIT}. \code{storage} holds
\code{itemSize * capacity} bytes and stays valid for as long as the queue is used.
\apichecked{a NULL storage, an item size or a capacity of 0.}

\apihead{rtos\_queueSend, rtos\_queueSendFromIsr}
\begin{lstlisting}
bool rtos_queueSend(rtos_queue_t *pQ, const void *pItem, rtos_tick_t timeout);
bool rtos_queueSendFromIsr(rtos_queue_t *pQ, const void *pItem);
\end{lstlisting}
Copies an item into the queue, waiting up to \code{timeout} for room (the ISR variant
never waits). If a task is blocked in receive, the item goes straight into that task's
buffer. While the caller waits, the item is read at the moment a receiver frees a slot:
it must not change until the call returns. A failed send is counted as dropped.
\apireturns{\code{true} if the item was queued.}

\apihead{rtos\_queueReceive, rtos\_queueReceiveFromIsr}
\begin{lstlisting}
bool rtos_queueReceive(rtos_queue_t *pQ, void *pItem, rtos_tick_t timeout);
bool rtos_queueReceiveFromIsr(rtos_queue_t *pQ, void *pItem);
\end{lstlisting}
Copies the oldest item out, waiting up to \code{timeout} for one. If a sender is blocked
because the queue was full, its item moves into the freed slot.
\apireturns{\code{true} if an item was received.}

\apihead{rtos\_queueCount}
\begin{lstlisting}
uint8_t rtos_queueCount(const rtos_queue_t *pQ);
\end{lstlisting}
\apireturns{the items stored now; a snapshot.} \apicontext{any.}

% -----------------------------------------------------------------------------------------
\section{Mailboxes}

\apihead{rtos\_mboxInit, rtos\_mboxPost, rtos\_mboxPostFromIsr, rtos\_mboxFetch, rtos\_mboxIsFull}
\begin{lstlisting}
void rtos_mboxInit(rtos_mbox_t *pMb, uint8_t *storage, uint8_t itemSize);
bool rtos_mboxPost(rtos_mbox_t *pMb, const void *pItem);
bool rtos_mboxPostFromIsr(rtos_mbox_t *pMb, const void *pItem);
bool rtos_mboxFetch(rtos_mbox_t *pMb, void *pItem, rtos_tick_t timeout);
bool rtos_mboxIsFull(const rtos_mbox_t *pMb);
\end{lstlisting}
A post never blocks: it replaces an unread item (counted in \code{overwrites}, and the
post returns \code{true}) or, if a task is blocked in fetch, copies the item straight
into that task's buffer. A fetch takes the item and leaves the mailbox empty; it returns
\code{false} on timeout. \code{rtos\_mboxIsFull} tells whether an unread item waits.
\apichecked{\code{rtos\_mboxInit}: a NULL storage or an item size of 0.}

% -----------------------------------------------------------------------------------------
\section{Task notifications}

\apihead{rtos\_notify, rtos\_notifyFromIsr}
\begin{lstlisting}
void rtos_notify(uint8_t taskId, uint16_t value, uint8_t op);
void rtos_notifyFromIsr(uint8_t taskId, uint16_t value, uint8_t op);
\end{lstlisting}
Updates the task's 16 bit notification value and wakes the task if it waits in
\code{rtos\_notifyWait}, or in \code{rtos\_notifyTake} and the value is now non-zero.
\begin{apiparams}
\item[op] \code{RTOS\_NOTIFY\_SET} replaces the value, \code{RTOS\_NOTIFY\_OR} ORs
  \code{value} in, \code{RTOS\_NOTIFY\_INC} adds one, saturating at 0xFFFF.
\end{apiparams}
\apichecked{\code{taskId}.}

\apihead{rtos\_notifyWait}
\begin{lstlisting}
bool rtos_notifyWait(uint16_t *pValue, rtos_tick_t timeout);
\end{lstlisting}
Returns as soon as any notification has arrived since the last successful wait, even one
that set the value to 0, stores the value in \code{*pValue} (may be NULL) and resets it to 0.
\apicontext{task; the idle task may only poll.}

\apihead{rtos\_notifyTake}
\begin{lstlisting}
bool rtos_notifyTake(rtos_tick_t timeout);
\end{lstlisting}
Waits until the value is non-zero and decrements it by one: with
\code{RTOS\_NOTIFY\_INC} on the sending side, a counting semaphore without an object. A
task should use either \code{rtos\_notifyWait} or \code{rtos\_notifyTake}, not both.

% -----------------------------------------------------------------------------------------
\section{Software timers}

\apihead{rtos\_timerInit}
\begin{lstlisting}
void rtos_timerInit(rtos_timer_t *pT, rtos_timerFunc_t callback, void *arg);
\end{lstlisting}
Prepares an inactive timer. Where the callback runs depends on
\code{RTOS\_CFG\_TIMER\_TASK\_PRIO}: in the tick interrupt (0; only ISR body functions
allowed) or in the timer task.
\apichecked{a NULL callback.}

\apihead{rtos\_timerStart}
\begin{lstlisting}
void rtos_timerStart(rtos_timer_t *pT, rtos_tick_t delay, rtos_tick_t period);
\end{lstlisting}
Arms the timer, or restarts it if active. The first expiry is on the \code{delay}-th tick
from now (after \code{delay - 1} to \code{delay} periods; 0 and 1 both mean the next
tick); after that every \code{period} ticks exactly, counted from the previous expiry.
\code{period} 0 is a one shot timer. The time spent with interrupts disabled grows with
the number of active timers.
\apicontext{any.}
\apichecked{the timer was initialised; \code{delay} and \code{period} at most the limit.}

\apihead{rtos\_timerStop, rtos\_timerIsActive, rtos\_timerRemaining}
\begin{lstlisting}
void        rtos_timerStop(rtos_timer_t *pT);
bool        rtos_timerIsActive(const rtos_timer_t *pT);
rtos_tick_t rtos_timerRemaining(const rtos_timer_t *pT);
\end{lstlisting}
Disarms the timer (nothing happens if it is not active; with a timer task an expiry
still waiting for the task is dropped, a running callback finishes); tells whether it is
armed; returns the ticks to the next expiry (0 if inactive or due).
\apicontext{any.}

% -----------------------------------------------------------------------------------------
\section{Memory pools}

\apihead{rtos\_poolInit}
\begin{lstlisting}
void rtos_poolInit(rtos_pool_t *pPool, uint8_t *storage, uint16_t blockSize, uint8_t numBlocks);
\end{lstlisting}
Builds the free list over the storage; also required after \code{RTOS\_POOL\_INIT}. A
block is returned at the alignment of its offset in the storage, so storage meant for
aligned objects must be aligned itself and have a block size that is a multiple of that
alignment.
\apichecked{a NULL storage, a block size below \code{sizeof(void *)}, 0 blocks.}

\apihead{rtos\_poolAlloc, rtos\_poolAllocFromIsr}
\begin{lstlisting}
void *rtos_poolAlloc(rtos_pool_t *pPool, rtos_tick_t timeout);
void *rtos_poolAllocFromIsr(rtos_pool_t *pPool);
\end{lstlisting}
Gets a block, waiting up to \code{timeout} for one to be freed (the ISR variant never
waits). A failed allocation is counted.
\apireturns{the block, or NULL.}

\apihead{rtos\_poolFree, rtos\_poolFreeFromIsr, rtos\_poolFreeCount}
\begin{lstlisting}
void    rtos_poolFree(rtos_pool_t *pPool, void *pBlock);
void    rtos_poolFreeFromIsr(rtos_pool_t *pPool, void *pBlock);
uint8_t rtos_poolFreeCount(const rtos_pool_t *pPool);
\end{lstlisting}
Returns a block; a task waiting for one gets it directly.
\apichecked{a pointer outside the pool or not at the start of a block. A double free is
caught only when the pool would hold more free blocks than it has.}

% -----------------------------------------------------------------------------------------
\section{Event groups}

\apihead{rtos\_evgroupInit, rtos\_evgroupSet, rtos\_evgroupSetFromIsr, rtos\_evgroupClear, rtos\_evgroupGet}
\begin{lstlisting}
void       rtos_evgroupInit(rtos_evgroup_t *pG, rtos_evt_t initialBits);
void       rtos_evgroupSet(rtos_evgroup_t *pG, rtos_evt_t bits);
void       rtos_evgroupSetFromIsr(rtos_evgroup_t *pG, rtos_evt_t bits);
void       rtos_evgroupClear(rtos_evgroup_t *pG, rtos_evt_t bits);
rtos_evt_t rtos_evgroupGet(const rtos_evgroup_t *pG);
\end{lstlisting}
Set wakes every waiter whose condition is now met, highest priority first; a waiter that
asked for \code{RTOS\_WAIT\_CLEAR} takes its bits out as it is woken. Clear wakes nobody.
With waiters present, set checks every task slot with interrupts disabled.

\apihead{rtos\_evgroupWait}
\begin{lstlisting}
rtos_evt_t rtos_evgroupWait(rtos_evgroup_t *pG, rtos_evt_t mask, uint8_t options,
                            rtos_tick_t timeout);
\end{lstlisting}
Waits until bits of \code{mask} are set. \code{options} is \code{RTOS\_WAIT\_ANY} or
\code{RTOS\_WAIT\_ALL}, optionally with \code{RTOS\_WAIT\_CLEAR}.
\apireturns{the satisfying bits (with ALL: the mask); 0 on timeout.}
\apiasserted{\code{mask} not 0.}

% -----------------------------------------------------------------------------------------
\section{Suspend and resume}

\apihead{rtos\_suspend, rtos\_resume, rtos\_resumeFromIsr}
\begin{lstlisting}
void rtos_suspend(uint8_t taskId);
void rtos_resume(uint8_t taskId);
void rtos_resumeFromIsr(uint8_t taskId);
\end{lstlisting}
Suspend keeps the task from running until a resume; a blocked task stays in its wait, a
wake-up that arrives meanwhile takes effect at the resume. No nesting count; a task may
suspend itself. A suspended task keeps its mutexes. Resume does nothing to a task that is
not suspended.
\apichecked{\code{taskId}; not the idle task.}

% -----------------------------------------------------------------------------------------
\section{Heap}

\apihead{rtos\_heapInit, rtos\_malloc, rtos\_free, rtos\_heapStats}
\begin{lstlisting}
void  rtos_heapInit(void *area, size_t size);
void *rtos_malloc(size_t size);
void  rtos_free(void *p);
void  rtos_heapStats(rtos_heapStats_t *pStats);
\end{lstlisting}
A first fit allocator with coalescing over the area given once to \code{rtos\_heapInit}.
Blocks are aligned to \code{2 * sizeof(void *)}. The search runs with interrupts
disabled and its time grows with the number of free blocks. \code{rtos\_free} ignores
NULL.
\apicontext{task.}
\apichecked{a NULL or too small area; \code{rtos\_malloc} before \code{rtos\_heapInit};
a freed pointer outside the heap or already free. A pointer into the middle of a block is
not detected.}

% -----------------------------------------------------------------------------------------
\section{Diagnostics}

\apihead{rtos\_getCurrentTask}
\begin{lstlisting}
uint8_t rtos_getCurrentTask(void);
\end{lstlisting}
\apireturns{the running task's id; 0 in the idle task; in an ISR body the interrupted task.}

\apihead{rtos\_getOverrunCount}
\begin{lstlisting}
uint8_t rtos_getOverrunCount(uint8_t taskId, bool reset);
\end{lstlisting}
\apireturns{the period overruns of the task since the start or the last reset; saturates
at 255. \code{reset} clears the counter after reading.} \apichecked{\code{taskId} 0 to
\code{RTOS\_CFG\_NUM\_TASKS}.}

\apihead{rtos\_getStackFree}
\begin{lstlisting}
size_t rtos_getStackFree(uint8_t taskId);
\end{lstlisting}
\apireturns{the bytes at the bottom of the task's stack never used since creation; 0 for
the idle task, a task never created, or without \code{RTOS\_CFG\_STACK\_CHECK}. The scan
reads the stack byte by byte with interrupts enabled.}

\apihead{rtos\_getTimeUs}
\begin{lstlisting}
uint32_t rtos_getTimeUs(void);
\end{lstlisting}
\apireturns{microseconds since start from the port's free running time base; wraps after
71.6 minutes.} \apicontext{any.}

\apihead{rtos\_getTaskRunTime, rtos\_getCpuLoad}
\begin{lstlisting}
rtos_runtime_t rtos_getTaskRunTime(uint8_t taskId);
uint8_t        rtos_getCpuLoad(void);
\end{lstlisting}
With \code{RTOS\_CFG\_RUNTIME\_STATS}: the accumulated run time of a task in counts of
\code{RTOS\_RUNTIME\_HZ} (the running task's includes the current slice), and the share
of time not spent in the idle task since the previous call, 0 to 100. Every call of
\code{rtos\_getCpuLoad} starts a new interval, so one task should own the readings.

\apihead{rtos\_getTaskInfo}
\begin{lstlisting}
bool rtos_getTaskInfo(uint8_t taskId, rtos_taskInfo_t *pInfo);
\end{lstlisting}
Fills a snapshot: state (\code{RTOS\_TASK\_*}), priority, effective priority, wait kind
and object, timeout, stack size and free, run time, overruns.
\apireturns{\code{true}; \code{false} for an id out of range or a NULL \code{pInfo}
(not a panic).}

% -----------------------------------------------------------------------------------------
\section{Panic and reset cause}

\apihead{rtos\_panic}
\begin{lstlisting}
void rtos_panic(uint8_t reason, uint32_t arg1, uint32_t arg2) __attribute__((noreturn));
\end{lstlisting}
Disables interrupts, records the reason and the arguments in memory that survives a
reset, calls \code{rtos\_panicHook} and halts or resets (\code{RTOS\_CFG\_PANIC\_ACTION}).
Applications use \code{RTOS\_PANIC\_APPLICATION}.
\apicontext{any.}

\apihead{rtos\_getLastPanic, rtos\_getResetCause}
\begin{lstlisting}
bool    rtos_getLastPanic(rtos_panicInfo_t *pInfo);
uint8_t rtos_getResetCause(void);
\end{lstlisting}
The record left before the last reset, reported in one boot only (\code{pInfo} may be
NULL to ask only), and the cause of the last reset as the hardware reports it
(\code{RTOS\_RESET\_*}).

% -----------------------------------------------------------------------------------------
\section{Hooks}

\apihead{rtos\_idleHook, rtos\_panicHook, rtos\_stackWarnHook, rtos\_assertFailed, rtos\_traceHook}
\begin{lstlisting}
void rtos_idleHook(void);
void rtos_panicHook(const rtos_panicInfo_t *pInfo);
void rtos_stackWarnHook(uint8_t task, size_t stackFree);
void rtos_assertFailed(const char *file, uint16_t line) __attribute__((noreturn, cold));
void rtos_traceHook(uint8_t event, uint32_t a, uint32_t b);
\end{lstlisting}
Weak functions the application overrides. The idle hook runs in every idle loop and
must not block. The panic hook runs with interrupts disabled before the halt or reset and
must not call the kernel. The stack warning hook is called once per task whose margin is
gone. \code{rtos\_assertFailed} is called for a failed assertion; its default stores file
and line and panics, and an override must not return either. The trace hook is called
with interrupts disabled at every traced event (\code{RTOS\_TRACE\_SWITCH}, \code{BLOCK},
\code{READY}, \code{TIMER}, \code{PANIC}) and must not call the kernel.

% -----------------------------------------------------------------------------------------
\section{Watchdog}

\apihead{rtos\_wdtInit, rtos\_wdtMonitor, rtos\_wdtCheckIn}
\begin{lstlisting}
void rtos_wdtInit(uint32_t timeoutMs);
void rtos_wdtMonitor(uint8_t task, rtos_tick_t maxInterval);
void rtos_wdtCheckIn(void);
\end{lstlisting}
Starts the hardware watchdog (the timeout is rounded up to what the hardware offers;
the idle task kicks it only after \code{rtos\_start}); requires a task to check in at
least every \code{maxInterval} ticks (0 stops monitoring it; the first deadline is
\code{maxInterval} from now); tells the watchdog the calling task is alive.
\apichecked{\code{task}; \code{maxInterval} at most the limit.}
\apiasserted{\code{timeoutMs} above 0.}
